\section{Privacy}
\label{sec:privacy}

\Zoo{} 3.0 supports two complementary privacy paths via \Lux{}
chain topology (\textsc{lp-134}): \FHE{} via F-Chain
(\textsc{lp-013}) for confidential compute on encrypted data, and
\ZKP{} via Z-Chain (\textsc{lp-063}) for selective disclosure of
private facts. They are not alternatives; they solve different
problems.

\subsection{F-Chain \FHE{} (LP-013)}

F-Chain (FVM) is a fully-homomorphic encryption execution
environment. Contracts on F-Chain operate on ciphertexts: the chain
never sees the plaintext, and yet the result of a computation is
verifiable on-chain.

\paragraph{Confidential ERC-20 (LP-067).}
The Confidential ERC-20 standard runs balances and allowances as
ciphertexts. A confidential transfer ciphertext-decrypts the source
balance, ciphertext-subtracts the amount, ciphertext-adds to the
destination, and ciphertext-encrypts the new balances. Validators
verify the operation without ever seeing a plaintext balance. Useful
for:

\begin{itemize}[leftmargin=1.2em]
  \item Confidential trading of high-value collectibles where price
        movements should not be public until settlement.
  \item Per-\LLM{} token holdings that the holder wants to keep
        private for governance reasons.
  \item Treasury balances that the \DAO{} holds confidentially until
        formal disclosure.
\end{itemize}

\paragraph{Private inference.}
A user submits an encrypted prompt to F-Chain; an \FHE{}-aware Zen
inference worker performs the forward pass on the ciphertext; the
encrypted output is returned to the user, who decrypts locally. The
worker never sees the plaintext prompt or the plaintext output. This
is the canonical pattern for confidential AI in regulated workflows
(healthcare, securities research, internal counsel).

\paragraph{When to use F-Chain.}
\FHE{} is the right primitive when the property to protect is the
\emph{data}, the computation is non-trivial, and verifiers need
correctness without seeing the data. \FHE{} costs more than \ZKP{}
per operation but supports arbitrary functions on the encrypted
data.

\subsection{Z-Chain \ZKP{} (LP-063)}

Z-Chain (ZVM) is a zero-knowledge proof execution and registry
environment. Z-Chain runs a Groth16 / Halo2 / Plonky2 verifier as a
\GPU{} kernel and hosts a circuit registry. Contracts on \Zoo{} L1
dispatch verification jobs to Z-Chain via precompile.

\paragraph{Selective disclosure.}
A user proves ``I am over 18'' or ``I hold $\geq 1{,}000$ \$ZOO''
without revealing date of birth or actual balance. Identity
assertions live on A-Chain (AIVM, \textsc{lp-134}); proofs about
those assertions live on Z-Chain.

\paragraph{Rollup of identity assertions.}
For high-volume assertion flows --- e.g., per-block KYC checks for a
regulated trading venue --- Z-Chain rolls up batches of proofs
into a single Groth16 verification, amortising verification cost.

\paragraph{When to use Z-Chain.}
\ZKP{} is the right primitive when the property to protect is a
\emph{predicate over data}, the predicate is bounded (compileable
to a circuit), and verifiers need a yes/no answer with cryptographic
backing. \ZKP{} verification is much cheaper than \FHE{} computation
per operation but requires you to know the predicate in advance.

\subsection{Choosing between them}

\begin{longtable}{p{0.30\textwidth}p{0.30\textwidth}p{0.34\textwidth}}
\toprule
\textbf{Question} & \textbf{F-Chain (\FHE{})} & \textbf{Z-Chain (\ZKP{})} \\
\midrule
Hide values? & yes (whole computation) & no (only the witness) \\
Hide identity? & via key management & yes (predicate over identity) \\
Arbitrary compute? & yes & no (must be a circuit) \\
Per-op cost & high & low to moderate \\
Setup cost & key generation & circuit compile + trusted setup (Groth16) \\
\bottomrule
\end{longtable}

The \Zoo{} 3.0 wallet (\S\ref{sec:wallet-desktop}) exposes a privacy
mode toggle that picks the right chain for the user's intent: ``hide
this transfer'' routes to F-Chain Confidential ERC-20; ``prove
eligibility'' routes to Z-Chain selective disclosure.

\subsection{\GPU{}-native privacy}

Both privacy paths are \GPU{}-native end-to-end. F-Chain's TFHE
bootstrapping runs on the \GPU{} per LP-013. Z-Chain's Groth16
verification (and the prover for circuits compiled in advance) runs
on the \GPU{}. The privacy primitives in \Zoo{} 3.0 are not bolted on
at the application layer; they execute on the same substrate as the
rest of the chain.
